\documentclass{article}
\usepackage[T1]{fontenc}
\usepackage{lmodern}
\usepackage{graphicx}
\usepackage{amsmath,amssymb}
\usepackage[a4paper,margin=2cm]{geometry}
\usepackage[absolute,overlay]{textpos}
\setlength{\TPHorizModule}{1mm}
\setlength{\TPVertModule}{1mm}
\usepackage[indonesian]{babel}
\usepackage{hyperref}
\setlength{\emergencystretch}{2em}
% unit: Halaman 54

\begin{document}

% === Halaman 54 (python/native) ===

54

• Contoh:

 DYDUXRMHTVDVNQDQNWDYDUXRMHARTJGWNQD

 Kriptogram yang berulang: DYUDUXRMH dan NQD. 

 Jarak antara dua buah perulangan DYUDUXRMH = 18. 

 Semua faktor pembagi 18 : \{18, 9, 6, 3, 2\}  

 Jarak antara dua buah perulangan NQD =20. 

 Semua faktor pembagi 20 : \{20, 10, 5, 4, 2\}. 

 Irisan dari kedua buah himpunan tersebut adalah 2 

 Panjang kunci kemungkinan besar adalah 2.

\end{document}
